问题一涉及"鱼类资源量增加"、"鱼类增长的驱动机制"和"基础资源的承载压力变化"。若将四大家鱼简化为单一消费者功能群，并把水草、浮游植物、浮游动物和底栖饵料汇总成单一资源，模型虽仍能对上2025年资源量恢复至1.8倍的观测事实，但无法说明哪一层资源率先承压，也难以联系"水下荒漠"等底层生态退化现象。因此，为留下这部分解释能力，模型中必须分别表示四类基础资源与四大家鱼的基本食性对应关系。
